% Generated draft: requires mathematical and source review.
\documentclass[11pt]{article}
\usepackage[margin=25mm]{geometry}
\usepackage{fontspec}
\setmainfont{Latin Modern Roman}
\usepackage{amsmath,amssymb,mathtools,unicode-math}
\setmathfont{Latin Modern Math}
\usepackage{enumitem,needspace,xurl,hyperref}
\setlength{\parindent}{0pt}
\setlength{\parskip}{5pt}
\setlength{\emergencystretch}{3em}
\newcommand{\N}{\mathbb{N}}
\newcommand{\Z}{\mathbb{Z}}
\newcommand{\Q}{\mathbb{Q}}
\newcommand{\R}{\mathbb{R}}
\newcommand{\C}{\mathbb{C}}
\newcommand{\F}{\mathbb{F}}
\newcommand{\A}{\mathbb{A}}
\newcommand{\PP}{\mathbb{P}}
\newcommand{\E}{\mathbb{E}}
\newcommand{\eps}{\varepsilon}
\newcommand{\dd}{\,\mathrm{d}}
\begin{document}
\section*{UnsolvedMath ID: 700920}
\textbf{Source problem code:} catalogue entry
\subsection{Definitions and mathematical statement}
A variety of groups is a class of groups defined by identities in the group language. It is finitely based if there is a finite set of identities whose group models are exactly the members of the variety. For a variety \(\mathfrak V\) and a positive integer \(n\), let \(F_n(\mathfrak V)\) denote the relatively free group in \(\mathfrak V\) on the free generating system \(x_1,\ldots,x_n\): every assignment of these generators to elements of any group in \(\mathfrak V\) extends uniquely to a homomorphism. Interpret \(F_\infty(\mathfrak V)\) as the relatively free group on the countably infinite free generating system \((x_i)_{i\geq 1}\), with generator indices effectively encoded. The supplied record does not explicitly specify the cardinality intended by the subscript \(\infty\); this interpretation requires confirmation.

The word problem with respect to one of these generating systems is solvable if an algorithm, given any finite word in the generators and their inverses, decides whether that word represents the identity element. It is unsolvable if no such algorithm exists.

Does there exist a finitely based variety of groups \(\mathfrak V\) such that, for every positive integer \(n\), the word problem in \(F_n(\mathfrak V)\) is solvable, but the word problem in \(F_\infty(\mathfrak V)\), with respect to its free generating system, is unsolvable? The finite-rank requirement has quantifier order \(\forall n\geq 1\;\exists A_n\), where \(A_n\) is a decision algorithm for rank \(n\); a single algorithm uniform in \(n\) is not required by the supplied statement.
\subsection{Short English statement}
The question asks for a group variety defined by finitely many identities whose relatively free groups have decidable word problems at every finite positive rank, while its relatively free group of infinite rank has an undecidable word problem with respect to a free generating system. Countably infinite rank is the interpretation used here, subject to confirmation.
\subsection{Sources}
\begin{itemize}
\item Local UnsolvedMath registry, permanent ID \(700920\), source problem code \texttt{catalogue entry}, titled \emph{Kourovka Notebook Problem 13.2}. The supplied record labels the problem \emph{open}; this is supplied registry status, not an independently verified mathematical or literature assessment.
\item \emph{The Kourovka Notebook (Unsolved Problems in Group Theory)}, original problem 13.2, 13th issue (1995); attributed to M. I. Anokhin. The supplied background identifies the 21st issue edition (2026), arXiv:1401.0300v48. The supplied collection description names E. I. Khukhro and V. D. Mazurov as editors. No external source was independently consulted.
\end{itemize}
\end{document}
